\documentclass{article}

\usepackage{icml2026}
\usepackage{amsmath,amssymb,amsthm}
\usepackage{booktabs,graphicx,xcolor,microtype}
\usepackage{hyperref}
\usepackage{cleveref}
\usepackage{url}
\usepackage{multirow}
\usepackage{wrapfig}
\usepackage{enumitem}
\usepackage{tikz}
\usetikzlibrary{shapes,arrows,arrows.meta,positioning}
\usepackage{algorithm}
\usepackage{algorithmic}
\usepackage{pgfplots}
\usepgfplotslibrary{fillbetween}
\usepgfplotslibrary{groupplots}
\pgfplotsset{compat=1.18}
\graphicspath{{../figures/}}

\input{math_commands.tex}

\newcommand{\N}{\mathbb{N}}
\newcommand{\E}{\mathbb{E}}
\newcommand{\diag}{\mathrm{diag}}

\icmltitlerunning{Does the Optimal Hallucination Detector Depend on Model Scale?}

\begin{document}

\twocolumn[
\icmltitle{Does the Optimal Hallucination Detector for LLM Tool Calls\\
  Depend on Model Scale?}

\begin{icmlauthorlist}
\icmlauthor{Anonymous}{}
\end{icmlauthorlist}

%% De-anonymize after acceptance:
\icmlaffiliation{devoteam}{Devoteam, Paris}
\icmlaffiliation{aws}{Amazon Web Services}
\icmlaffiliation{openai}{OpenAI}
\icmlcorrespondingauthor{Valentin No\"{e}l}{val.noel@proton.me}

\icmlkeywords{agentic AI, tool-use hallucination, safety, spectral graph theory, interpretability}

\vskip 0.3in
]

\printAffiliationsAndNotice{}

%══════════════════════════════════════════════════════════════════════════════
\begin{abstract}
%══════════════════════════════════════════════════════════════════════════════
When an LLM agent hallucinates a tool call, it executes, no hedge,
no uncertainty signal.  Standard detection approaches fail structurally:
token log-probabilities and sampling consensus score near chance
(AUC 0.51\,--\,0.62) because hallucinated calls are syntactically valid
and high-confidence.  Two independently developed supervised approaches
can detect these failures from a single forward pass: \emph{spectral
attention topology}, which tracks algebraic connectivity of per-layer
attention graphs without accessing hidden states, and \emph{token-role
hidden-state probing}, which supervises an MLP on structural positions
within the generated call.  The natural question, which one to
deploy?, has had no controlled answer.

We provide the first head-to-head comparison across seven models
(1B to 14B) on identical data, splits, and labels.  The answer is
non-monotone.  Our \textbf{Layerwise Multi-Metric} (LMM) probe reaches
AUC 0.957 at 3B using only attention weights, exceeding Hidden (0.947);
at 8B, hallucination signal diffuses broadly across 32 layers (mean
per-layer AUC 0.59) and trajectory aggregation (SpRich, AUC 0.888)
dominates; at 14B, deeper layer specialisation re-concentrates the
signal and LMM recovers its lead (0.917 vs.\ 0.882).  Architecture
matters independently: grouped-query attention (GQA) collapses spectral
trajectory statistics to near chance (0.432) while per-layer LMM
remains robust.  Merging the two signals never helps.

The practical implication is precise: match the guardrail to model
depth and attention architecture, not to signal combination.
\end{abstract}

%══════════════════════════════════════════════════════════════════════════════
\section{Introduction}
\label{sec:intro}
%══════════════════════════════════════════════════════════════════════════════

LLM agents that invoke external tools (APIs, databases, privileged
actions) fail silently when they hallucinate a tool call.  Unlike
open-ended generation, the output is syntactically valid and
high-confidence, a plausible function name and argument schema that
happens to be wrong.  Downstream pipeline steps consume it without
question, so a malformed transaction or misconfigured
infrastructure command can execute before any human reviews it.

Detection must therefore happen at generation time, from a single
forward pass, with no ground-truth oracle and no extra sampling budget.
Standard uncertainty signals fail structurally on this task: token
log-probabilities reach AUC 0.51\,--\,0.62 and $N{=}5$ sampling
consensus is no better at $5\times$ the cost
(Table~\ref{tab:baselines}).  The reason is not noise, a model does
not express uncertainty when hallucinating a tool call it is confident
about.

Two independent lines of work have shown that supervised signals from
the same forward pass can break through this barrier.  Spectral
attention-graph diagnostics \cite{noel2026geometry} track algebraic
connectivity (Fiedler value $\lambda_2$) of per-layer graph Laplacians
with no hidden-state access.  A token-role hidden-state probe
\cite{healy2026internal} supervises an MLP on the function-name,
argument-value, and closing-delimiter positions of the generated call.
Both approaches read different tensors from the same pass and were
developed independently.  The question they leave open, \emph{which
one to deploy, and do they complement each other?}, is what this
paper answers.

Our three contributions:
\begin{enumerate}[noitemsep,topsep=2pt,leftmargin=*]
  \item A \textbf{controlled head-to-head} across seven models (1B to
    14B), identical splits and labels, including three concurrent
    spectral and hidden-state baselines
    \cite{binkowski2025hallucination,li2025llm,badash2026between}.
  \item A \textbf{Layerwise Multi-Metric} (LMM) probe that evaluates
    spectral features per layer, recovering signal that trajectory
    aggregation discards.
  \item A \textbf{precise deployment rule}: LMM dominates at
    ${\leq}$3B and ${\geq}$14B; trajectory aggregation (SpRich)
    dominates at 8B; GQA attention architecture
    collapses SpRich while leaving LMM robust; merging never helps.
\end{enumerate}

\begin{figure*}[t]
    \centering
    \begin{tikzpicture}[
    node distance=0.8cm and 1.2cm,
    block/.style={
        rectangle, draw, fill=white,
        inner sep=3pt, align=center,
        rounded corners=3pt, font=\small\sffamily,
        line width=0.6pt
    },
    smallblock/.style={
        rectangle, draw, fill=white,
        inner sep=4pt, align=center,
        rounded corners=2pt, font=\scriptsize\sffamily,
        line width=0.5pt
    },
    arrow/.style={-{Stealth[scale=1.1]}, thick, draw=gray!80},
    ]

    \node (input) [block, fill=gray!10, text width=2.2cm] {
        \textbf{Tool Call}\\[2pt]
        {\tiny\ttfamily transfer\_funds(}\\
        {\tiny\ttfamily \ \ acc="1234",}\\
        {\tiny\ttfamily \ \ amt=500)}
    };

    \node (forward) [block, right=of input, fill=orange!10, text width=1.8cm] {
        \textbf{Single}\\
        \textbf{Forward Pass}
    };

    \node (attention) [smallblock, fill=blue!5, draw=blue!40, above right=0.3cm and 1cm of forward, text width=2.4cm] {
        \textbf{Attention Routing}\\[2pt]
        {\tiny Graph Laplacian $\mathbf{L}_\ell$}
    };

    \node (trajectory) [smallblock, fill=blue!10, draw=blue!60, right=of attention, text width=2.6cm] {
        \textbf{Spectral Trajectory}\\[2pt]
        \begin{tikzpicture}[scale=0.6]
            \draw[->, very thin, gray] (0,0) -- (2.2,0) node[right, font=\tiny] {$\ell$};
            \draw[->, very thin, gray] (0,0) -- (0,1) node[above, font=\tiny] {$\lambda_2$};
            \draw[blue!60, thick] (0,0.8) -- (0.5,0.75) -- (1.0,0.6) -- (1.5,0.7) -- (2.0,0.8);
            \draw[red!60, thick] (0,0.8) -- (0.5,0.7) -- (1.0,0.2) -- (1.5,0.15) -- (2.0,0.1);
        \end{tikzpicture}\\[2pt]
    };

    \node (hidden) [smallblock, fill=green!5, draw=green!40, below right=0.2cm and 1cm of forward, text width=2.4cm] {
        \textbf{Hidden States}\\[2pt]
        {\tiny Token-role $\mathbf{h}_\ell$}
    };

    \node (probe) [smallblock, fill=green!10, draw=green!60, right=of hidden, text width=3.6cm] {
        \textbf{Token-Role Probe}\\[4pt]
        {\tiny $z_\ell = h[t_\text{func}] \,\|\, \bar{h}[T_\text{args}] \,\|\, h[t_\text{end}]$}\\[4pt]
    };

    \node (merger) [block, fill=purple!10, draw=purple!50, below = 1.6cm of trajectory, right= 1cm of forward, text width=3.4cm] {
        \textbf{Convex Merger}\\[4pt]
        {\scriptsize $s = \alpha\, s_\text{sp} + (1{-}\alpha)\, s_\text{hid}$}
    };

    \node (output) [block, fill=gray!15, right=of merger, text width=2.8cm] {
        \textbf{Guardrail}\\
    };

    \draw[arrow] (input) -- (forward);
    \draw[arrow] (forward.east) -- ++(0.3,0) |- (attention.west);
    \draw[arrow] (forward.east) -- ++(0.3,0) |- (hidden.west);
    \draw[arrow] (attention) -- (trajectory);
    \draw[arrow] (trajectory.east) -- ++(1.25,0) |- (merger.9);
    \draw[arrow] (hidden) -- (probe);
    \draw[arrow] (probe.east) -- ++(0.3,0) |- (merger.351);
    \draw[arrow] (merger) -- (output);

    \node[above=0.1cm of trajectory, font=\tiny\itshape, text=gray!80] {SpRich / LMM};
    \node[below=0.1cm of probe, font=\tiny\itshape, text=gray!80] {Hidden \cite{healy2026internal}};

    \end{tikzpicture}
    \caption{\textbf{Architecture.}
      Two independent detection signals from a single forward pass.
      Top: spectral topology of attention routing (Fiedler value per layer).
      Bottom: token-role hidden-state probe \cite{healy2026internal}.}
    \label{fig:architecture}
\end{figure*}


%══════════════════════════════════════════════════════════════════════════════
\section{Related Work}
\label{sec:related}
%══════════════════════════════════════════════════════════════════════════════

\paragraph{Tool-use benchmarks.}
Gorilla \cite{patil2024gorilla} and ToolLLM \cite{qin2023toolllm}
benchmark tool-calling capability but not the internal failure modes
that produce confident wrong outputs.  The token-role hidden-state
probe \cite{healy2026internal} was the first to demonstrate that
supervised probes on hidden states can detect tool-use hallucination
from a single forward pass, establishing a representation built from the
function name, argument mean, and closing delimiter as a reliable input
to lightweight classifiers.  That probe forms one of the two signals
we evaluate; we extend it to a multi-layer setting and provide the
first comparison against an attention-only alternative.

\paragraph{Internal-state hallucination detection.}
INSIDE \cite{chen2024inside} uses hidden-state covariance; HALT
\cite{bhatnagar2026halt} extends probing to factual claims; process
reward models \cite{lightman2023lets} verify reasoning steps.  All
require hidden-state access.  Our spectral signal operates on attention
matrices alone.

\paragraph{Spectral methods for transformers.}
Attention-graph Laplacians were introduced for hallucination detection
in \cite{noel2025gsp}.  Concurrent work confirms the utility of
spectral geometry: LapEigvals \cite{binkowski2025hallucination} uses
top-$k$ eigenvalues; HSAD \cite{li2025llm} applies FFT to hidden-layer
dynamics; Cross-Layer Agreement \cite{badash2026between} measures
attention consistency across depth.  We evaluate all three as baselines.


%══════════════════════════════════════════════════════════════════════════════
\section{Method}
\label{sec:method}
%══════════════════════════════════════════════════════════════════════════════

\paragraph{Task and data.}
We sample $N{=}750$--$1000$ examples per model from the Glaive
function-calling benchmark \cite{glaive2024function} across seven
models: Llama-3.2-1B, Gemma-4-2B, Llama-3.2-3B, Qwen-2.5-3B,
Phi-4-Mini (3.8B), Llama-3.1-8B, and Phi-4 (14B, 4-bit NF4).  Labels
are assigned by JSON schema validation against the ground-truth call.  A
template-consistent 70/15/15 partition (SHA-256 hash of the ground-truth
string) ensures no prompt template seen at training appears at test
time, matching the protocol of \cite{healy2026internal}.

\paragraph{Signal 1: Spectral attention topology.}
For each layer~$\ell$, we average multi-head attention matrices and
symmetrize to form an adjacency $W_\ell$.  The graph Laplacian $L_\ell =
D_\ell - W_\ell$ yields the Fiedler value $\lambda_2(L_\ell)$, which
measures algebraic connectivity: low values indicate fragmented routing,
a pattern associated with hallucination \cite{noel2026geometry}.  We
extract five spectral metrics per layer (Fiedler value, energy,
smoothness index, spectral entropy, high-frequency energy ratio) and
represent each via 15 trajectory statistics (mean, slope, skewness,
dominant FFT harmonics, and others), giving a 135 to 200 dimensional
feature vector.  A logistic regression with $\ell_2$ regularization, $C$
tuned on validation, defines the \textbf{SpRich} probe.  For efficiency,
we use a subgraph-based Lanczos approximation (Fast-Fiedler) that
reduces per-sample spectral cost from 3.3\,s to 19\,ms.

This signal requires only attention weight matrices and no hidden-state
access, a critical advantage for API-served models where internal
activations are unavailable.

\paragraph{Signal 2: Hidden-state probe (Hidden).}
Following \cite{healy2026internal}, we extract token-role representations
at three structural positions within the generated call:
\begin{equation}
  z_\ell = h_\ell[t_\text{func}] \;\|\;
           \text{mean}(h_\ell[T_\text{args}]) \;\|\;
           h_\ell[t_\text{end}],
  \label{eq:tokenrole}
\end{equation}
where $t_\text{func}$ is the function-name token, $T_\text{args}$ are
argument-value positions, and $t_\text{end}$ is the closing delimiter.
We concatenate across 8 evenly spaced layers and train a
single-hidden-layer MLP (512 units, AdamW, cosine schedule, early
stopping).

\paragraph{Layerwise Multi-Metric (LMM) probe.}
Trajectory aggregation (SpRich) reduces $\lambda_2$ across all layers to
summary statistics, discarding information about which layers exhibit
anomalous routing.  LMM treats the raw $L \times 5$ matrix of per-layer
metric values as a direct feature vector.  We evaluate two variants:
\textbf{LMM-LR} (logistic regression) and \textbf{LMM-GBT}
(gradient-boosted trees, depth 3, 200 estimators).  Each GBT split
implements a learned threshold on a specific (layer, metric) pair.
No hidden-state access is required.

\paragraph{Convex merger.}
We test complementarity via $s = \alpha\, s_\text{sp} + (1{-}\alpha)\,
s_\text{hid}$, with $\alpha^*$ selected over 41 values by maximizing
validation AUC.


%══════════════════════════════════════════════════════════════════════════════
\section{Results}
\label{sec:results}
%══════════════════════════════════════════════════════════════════════════════

\subsection{Standard Signals Fail}

\begin{table}[t]
\centering
\caption{Unsupervised baselines vs.\ the best single spectral metric at
  a Youden-$J$ threshold (Sweep, no training).
  $\dagger$: teacher-forcing logprob (generation-based logprob was not run).
  $\ddagger$: reversed scoring direction---high logprob / high consensus
  predicts hallucination (the model commits confidently to wrong calls);
  standard direction gives 0.346 and 0.180 respectively.
  $\S$: $N{=}0$ sampling (14B inference cost); all consensus scores
  identical, AUC trivially 0.500.}
\label{tab:baselines}
\small
\begin{tabular}{lccc}
\toprule
Model & Logprobs & Consensus & Sweep \\
\midrule
Llama-1B          & 0.510 & 0.508 & \textbf{0.773} \\
Gemma 4 2B        & $0.672^{\dagger}$ & $0.820^{\dagger\ddagger}$ & \textbf{0.874} \\
Llama-3B          & 0.615 & 0.505 & \textbf{0.767} \\
Qwen-2.5-3B       & 0.579 & 0.512 & \textbf{0.818} \\
Phi-4 (14B)       & $0.655^{\dagger\ddagger}$ & $0.500^{\dagger}$ & \textbf{0.850} \\
\bottomrule
\end{tabular}
\end{table}

Token log-probabilities and $N{=}5$ sampling consensus are near chance
on every model where they were computed (Table \ref{tab:baselines}).
Even a single spectral metric with no training exceeds both by 15 to
26 AUC points, confirming that tool-use hallucination is invisible to
standard uncertainty signals but already detectable in the topology of
attention graphs.  Sweep peaks at 0.874 for Gemma 4 2B and 0.850 for
Phi-4 14B, the strongest unsupervised results across all seven models.

\subsection{Scale-Dependent Crossover}

\begin{table}[t]
\centering
\setlength{\tabcolsep}{1pt}
\caption{Test AUC with 95\% bootstrap CIs (1\,000 resamples).  Best per
  row in \textbf{bold}.  $\dagger$: Qwen uses LMM-LR (GBT overfits the
  reduced feature space of this hybrid architecture).
  $^\S$: Phi-4 Mini uses GQA (8 KV heads, 32 query heads): queries
  sharing a KV pair produce identical rows in the attention matrix,
  making the graph Laplacian rank-deficient. }
\label{tab:results}
\footnotesize
\begin{tabular}{lccc}
\toprule
Model & Hidden & LMM-GBT & SpRich \\
\midrule
Llama-1B          & \textbf{0.922}\,\tiny{[.873,.963]} & 0.918\,\tiny{[.857,.971]} & 0.856\,\tiny{[.784,.919]} \\
Gemma 4 2B        & \textbf{0.967}\,\tiny{[.935,.990]} & 0.918\,\tiny{[.858,.969]} & 0.894\,\tiny{[.827,.950]} \\
Llama-3B          & 0.947\,\tiny{[.908,.980]} & \textbf{0.957}\,\tiny{[.922,.985]} & 0.896\,\tiny{[.833,.944]} \\
Qwen-2.5-3B      & \textbf{0.948}\,\tiny{[.909,.982]} & 0.933$^\dagger$\,\tiny{[.885,.974]} & 0.905\,\tiny{[.854,.955]} \\
Phi-4 Mini (3.8B) & \textbf{0.969}\,\tiny{[.934,.999]} & 0.924\,\tiny{[.865,.972]} & 0.432$^\S$\,\tiny{[.327,.543]} \\
Llama-8B          & 0.848\,\tiny{[.778,.911]} & 0.870\,\tiny{[.811,.922]} & \textbf{0.888}\,\tiny{[.831,.935]} \\
Phi-4 (14B)       & 0.882\,\tiny{[.824,.933]} & \textbf{0.917}\,\tiny{[.873,.959]} & 0.907\,\tiny{[.853,.953]} \\
\bottomrule
\end{tabular}
\end{table}

Table \ref{tab:results} reveals a non-monotone pattern.  At
${\leq}$2B, Hidden dominates (0.922 at 1B, 0.967 at Gemma 4 2B);
LMM-GBT is close behind (0.918 both), but SpRich lags (0.856--0.894),
as trajectory aggregation loses layer-specific signal at low depth.

At 3B, LMM-GBT overtakes Hidden for the first time: 0.957 vs.\ 0.947,
using only attention weights.  The gain is recovered per-layer
resolution, SpRich (0.896) shows that trajectory aggregation
discards it.

Phi-4 Mini (3.8B) introduces an architectural confound.  Hidden (0.969)
and LMM-GBT (0.924) remain strong, but SpRich collapses to 0.432,
below chance.  The cause is GQA (8 KV heads, 32 query heads): queries
sharing a KV pair produce correlated attention distributions, yielding
near-singular Laplacians whose trajectory statistics carry no
discriminative signal.  LMM-GBT evaluates each layer's Fiedler value
independently and stays robust.

At 8B the crossover inverts: SpRich (0.888) beats LMM-GBT (0.870) and
Hidden (0.848).  Post-hoc analysis reveals why, train AUC 1.000,
test AUC 0.870, a 0.13-point gap absent at smaller scales.  Per-layer
Fiedler analysis shows 24 of 32 layers carry a signal (AUC${>}$0.55),
but each is weak (mean 0.59, peak 0.74, std 0.11).  At 1--3B,
hallucination creates sharp routing anomalies in a few layers GBT can
exploit; at 8B the signal is too diffuse for 160-feature GBT on 700
samples.  SpRich's global trajectory pooling handles diffuse signal
without per-layer fitting.

At 14B LMM-GBT recovers its lead (0.917 vs.\ SpRich 0.907, Hidden
0.882).  Forty layers provide 200 per-layer features with stronger
specialisation; the diffusion regime that hurt GBT at 32 layers
resolves as depth and per-layer signal variance both grow.

The pattern is clear: Hidden at ${\leq}$2B, LMM-GBT at 3B and
${\geq}$14B, SpRich at 8B, SpRich collapsed for GQA.  Detector choice
must account for model scale and attention architecture.

\subsection{Merging Does Not Help}

\begin{table}[t]
\centering
\setlength{\tabcolsep}{2pt}
\caption{Convex merger: $\alpha^*$ tuned on validation.  CIs overlap the
  stronger standalone method in every case.}
\label{tab:merger}
\footnotesize
\begin{tabular}{lcc}
\toprule
Model & Merger AUC & $\alpha^*$ \\
\midrule
Llama-1B          & 0.926\,\tiny{[.875,.964]} & 0.05 \\
Gemma 4 2B        & 0.967\,\tiny{[.933,.992]} & 0.10 \\
Llama-3B          & 0.947\,\tiny{[.900,.981]} & 0.00 \\
Qwen-2.5-3B      & 0.948\,\tiny{[.906,.982]} & 0.03 \\
Phi-4 Mini (3.8B) & 0.964\,\tiny{[.923,.996]} & 0.10 \\
Llama-8B          & 0.868\,\tiny{[.798,.927]} & 0.65 \\
Phi-4 (14B)       & 0.913\,\tiny{[.858,.956]} & 0.53 \\
\bottomrule
\end{tabular}
\end{table}

The merger CI overlaps the stronger standalone at every scale
(Table~\ref{tab:merger}).  At small scales $\alpha^*{\approx}0$ or
$0.10$, the merger inherits the Hidden score.  LMM-GBT catches more
exclusive positives than Hidden at every scale (14--31\% of
hallucinated test samples vs.\ 0--7\% Hidden-only), but these
exclusively caught samples coincide with noisy Hidden scores, so the
signals detect overlapping failure populations, not complementary ones.
More expressive fusion architectures likewise fail (see Appendix).
The conclusion is unambiguous: choose by scale and architecture; do not
combine.

\subsection{Concurrent Baselines}

\begin{table}[t]
\centering
\setlength{\tabcolsep}{2pt}
\caption{Comparison to concurrent spectral and hidden-state baselines.
  All evaluated on identical template-consistent splits.}
\label{tab:concurrent}
\footnotesize
\begin{tabular}{lcccc}
\toprule
Method & 1B & 3B & Q-3B & 8B \\
\midrule
HSAD \cite{li2025llm}           & 0.624 & 0.888 & 0.823 & 0.686 \\
Cross-Layer \cite{badash2026between} & 0.814 & 0.894 & 0.839 & 0.659 \\
LapEigvals \cite{binkowski2025hallucination} & 0.772 & 0.899 & 0.623 & 0.851 \\
\midrule
LMM-GBT [ours]  & 0.918          & \textbf{0.957} & 0.933 & 0.870 \\
SpRich [ours]    & 0.856          & 0.896          & 0.905 & \textbf{0.888} \\
Hidden [ours]    & \textbf{0.922} & 0.947          & \textbf{0.948} & 0.848 \\
\bottomrule
\end{tabular}
\end{table}

LMM-GBT outperforms all concurrent baselines at every scale
(Table \ref{tab:concurrent}).  HSAD and Cross-Layer show high
sensitivity to model scale (AUC ranges of 0.26 and 0.24 respectively).
At 8B, Hidden (0.848) trails LapEigvals (0.851) by 3 points, where
SpRich leads (0.888); at all other scales Hidden and LMM-GBT both
dominate.  LapEigvals, which uses raw eigenvalues without trajectory
statistics, degrades sharply on Qwen's hybrid attention architecture
(0.623), while SpRich remains robust (0.905), suggesting that
trajectory-level statistics provide the regularization needed for
architectural generalization.

\subsection{Cross-Benchmark Transfer (BFCL)}

\begin{table}[t]
\centering
\setlength{\tabcolsep}{2pt}
\caption{BFCL v3 AST hard-mode benchmark \cite{patil2025bfcl} ($N{=}500$, exclusively from \texttt{parallel}, \texttt{multiple}, and \texttt{parallel\_multiple} categories.}
\label{tab:bfcl}
\footnotesize
\begin{tabular}{lccc}
\toprule
Model & LMM-GBT & Hidden & Merger \\
\midrule
Llama-3B     & 0.894\,\tiny{[.764,.986]} & \textbf{0.898}\,\tiny{[.750,1.0]} & \textbf{0.898}\,\tiny{[.759,1.0]} \\
Qwen-2.5-3B & 0.742\,\tiny{[.565,.896]} & 0.800\,\tiny{[.626,.945]} & \textbf{0.838}\,\tiny{[.693,.952]} \\
Llama-8B     & \textbf{0.842}\,\tiny{[.719,.942]} & 0.679\,\tiny{[.538,.828]} & \textbf{0.842}\,\tiny{[.716,.944]} \\
\bottomrule
\end{tabular}
\end{table}

The scale-dependent pattern replicates on BFCL v3 hard mode
(Table \ref{tab:bfcl}): at 3B, Hidden and LMM-GBT tie; at 8B,
LMM-GBT leads by 16 AUC points.  Every BFCL sample requires detecting
hallucinations across 2--4 concurrent tool calls, yet detectors trained
on single calls transfer without modification, suggesting the
spectral hallucination signature is call-local rather than
context-global, which supports per-step deployment in tool chains.

%──────────────────────────────────────────────────────────────────────────────
\section{Discussion}
\label{sec:discussion}
%──────────────────────────────────────────────────────────────────────────────

\paragraph{Two signals, two mechanisms.}
The spectral signal captures attention routing: hallucinated calls fragment the attention graph, decreasing the Fiedler value. The hidden-state signal reflects representations at structural tokens. These mechanisms are correlated rather than complementary, detecting overlapping failures; thus, merging offers no gain. At ${\leq}$3B and ${\geq}$14B, LMM-GBT on attention matrices (AUC $\leq$ 0.957) matches or exceeds Hidden-state performance, suiting API-only access. At 8B, where hallucinations diffuse across 24 layers, SpRich's global trajectory pooling outperforms per-layer GBT, though architecture independently collapses SpRich for GQA models.

\paragraph{Agentic deployment.}
Our detectors process single calls and require step-by-step application in chains. Table~\ref{tab:bfcl} suggests single-call training transfers to multi-call settings. However, true sequential chains---where call $t{+}1$ inherits hallucinations from call $t$---remain untested. This is a critical open problem: error rates will grow super-linearly if detection fails at any step.

\paragraph{Considerations and limitations.}
Fast-Fiedler reduces spectral overhead from 3.3\,s to 19\,ms via GPU Lanczos on small subgraphs. For zero-shot settings, LMM-Cascade (calibration-only) yields AUC 0.72--0.81. While we studied 1B--14B models, the non-monotone crossover suggests LMM-GBT may lead at 70B+ as layer specialization re-concentrates signals. Finally, learned fusion remains infeasible due to extreme dimensionality mismatch (${\approx}10^5$ vs. ${\approx}10^2$); larger benchmarks are needed to bridge this gap.

%══════════════════════════════════════════════════════════════════════════════
\section{Conclusion}
%══════════════════════════════════════════════════════════════════════════════

The answer to our title question is yes, and the dependence is
non-monotone.  At 3B, spectral LMM alone (AUC 0.957) matches or
exceeds Hidden (0.947) without accessing hidden states; at 8B, diffuse
hallucination signal reverses the advantage to SpRich (0.888); at 14B,
LMM recovers (0.917 vs.\ 0.882).  GQA architecture collapses spectral
trajectory statistics entirely while per-layer LMM stays robust.  The
practical rule: route to LMM-GBT at ${\leq}$3B and ${\geq}$14B, to
SpRich at 8B, and check attention architecture before choosing any
spectral detector.  


%══════════════════════════════════════════════════════════════════════════════
\newpage
\bibliography{references}
\bibliographystyle{icml2026}

\end{document}